\subsection{Cross-regime routing and physical-space T2-C assembly}
\label{sec:outer_routing}

Unlike the internal routing in Section~\ref{sec:hybrid_specialists},
which selects correction experts within a local specialist, the outer
E2 router produces a preference over the S/H/P specialists using only
the physical parameter:
\begin{equation}
  \widetilde\mu_R=(\mu-\mu_R)/\sigma_R,\qquad
  \bm\pi(\mu)=\operatorname{softmax}
  \bigl(\mathcal R_{E2}(\widetilde\mu_R)\bigr).
  \label{eq:main_e2_preference}
\end{equation}
For the verified centered-square router,
$\mathcal R_{E2}(x)=W_2\tanh(W_1x+\bm b_1)+\bm b_2$
is a $1$--$24$--$3$ MLP with 123 parameters.
It is trained on chart-assigned parameter labels, so $\bm\pi$ is a
parameter-conditioned preference, not an applicability or stability
probability. The deterministic parameter-only output is identical
for windows with the same $\mu$ and remains fixed during each rollout.
The E2 normalization and training protocol are specified in
Appendix~\ref{app:implementation}.

\paragraph{Prespecified adjacent candidates.}
In the reported Square overlap experiments, the evaluation protocol
supplies an ordered candidate pair $(r,s)=(S,H)$ or $(P,H)$ on
prespecified common chart support, with $K=24$ prediction steps.
E2 does not generate this candidate set, and no learned or
validation-thresholded applicability mask is used.
For the supplied pair, define
\begin{equation}
  \bar\pi_r=\frac{\pi_r}{\pi_r+\pi_s},\qquad
  \bar\pi_s=1-\bar\pi_r.
  \label{eq:main_pair_preference}
\end{equation}
E2 Top-1 selects $r$ when $\bar\pi_r\geq1/2$ and $s$ otherwise.
The pair order means that the first-candidate weight always refers
to Steady in the S--H study and Periodic in the H--P study.

\paragraph{History-conditioned output fusion.}
Both candidates encode the same physical initialization history in
their own charts and advance independently using their configured
local predictors. T2-C adjusts the E2 pair preference using a
chart-independent descriptor of the initial physical history:
\begin{equation}
  \widehat{\bm q}_{\rm T2C}(t)
  =\alpha\,\widehat{\bm q}^{(r)}(t)
   +(1-\alpha)\widehat{\bm q}^{(s)}(t),\qquad
  \alpha=\operatorname{sigmoid}
  \left(\operatorname{logit}(\bar\pi_r^{\rm clip})
       +c_\psi([\widetilde\mu_G;\bm d_n])\right).
  \label{eq:main_t2c}
\end{equation}
Here $\bar\pi_r^{\rm clip}$ is $\bar\pi_r$ clipped to
$[10^{-6},1-10^{-6}]$, and $\bm q$ collects reconstructed velocity
and gauge-corrected pressure on the common physical mesh.
The eight descriptor channels $\bm d_n$ are computed from the initial
three physical states, without future reference states.
The gate input $[\widetilde\mu_G;\bm d_n]$ uses training-window
normalization, distinct from the trajectory-label normalization
$\widetilde\mu_R$ used by E2. The correction network is a
$9$--$64$--$64$--$1$ MLP with SiLU hidden activations and a
zero-initialized final layer. Its construction and physical-field
training loss are given in Appendix~\ref{app:routing}.

The scalar $\alpha$ is computed for each prediction window and is
shared by velocity and pressure throughout that window.
Unlike E2, it may differ between histories at the same parameter.
Fusion acts only on reconstructed outputs and is never fed back
into either candidate. This interface does not require identical
internal architectures for the two local predictors.

\paragraph{Evaluation scope.}
These results evaluate conditional routing and assembly on prescribed
overlap support, not automatic applicability detection or
out-of-distribution rejection. The offline cache checks, including
their use of reference-field amplitudes, are disclosed in
Appendix~\ref{app:routing}; they are not deployable masks.
Output fusion preserves a common linear constraint only when both
candidates satisfy it under the same discrete operator, and it need
not preserve nonlinear momentum or PPE constraints.
Theoretical identities and their assumptions are given in
Appendix~\ref{app:fusion_consistency}; the present error tables do not
establish measured FVM residual or long-time phase consistency.

